\section{RW12 supplementary digital fabrication handoff}
\label{sec:rw12}

RW12 is included because it is practically useful, not because it strengthens
Theorem~\ref{thm:wrapper}.  It is a corrected body-preserving digital handoff
for \treeA{} containing fabrication-oriented files, checklist material, and a
compact package zip.  Its governing report is
\path{results/rw12_external_fabrication_review_package/rw12_external_fabrication_review_report.json}.

\begin{certprop}[RW12 hash and claim boundary]
The public checker verifies that the RW12 zip hash matches the RW12 report and
that RW12 does not promote a physical-validation claim.  The RW12 package is
therefore an appendix-level digital handoff, not theorem evidence.
\end{certprop}

The practical interpretation is straightforward: RW12 can be passed to a CAD,
fabrication, or external-review workflow as a digital artifact.  It does not
certify positive-thickness hinge clearances, material tolerances, assembly
success, printed prototype behavior, or measured motion.

\begin{remark}
This distinction matters for publication.  The mathematical result is the
zero-thickness one-parameter certificate package.  A real-world mechanism would
require additional positive-thickness modelling, tolerances, fabrication
measurements, and physical validation.  Those are future work, not hidden
assumptions of the present theorem.
\end{remark}
